\documentclass[11pt]{article}

\usepackage{amsmath,amssymb,amsfonts}
\usepackage{bm}
\usepackage{geometry}
\geometry{margin=1in}
\usepackage{hyperref}
\usepackage{booktabs}

\newcommand{\uu}{\bm{u}}
\newcommand{\xx}{\bm{x}}
\newcommand{\ff}{\bm{f}}
\newcommand{\nn}{\bm{n}}
\newcommand{\grad}{\nabla}
\newcommand{\dive}{\grad\!\cdot\!}
\newcommand{\eps}{\bm{\varepsilon}}
\newcommand{\Real}{\mathbb{R}}

\title{A Manufactured Solution with the \emph{Cocquet} Formulation:\\
       Verifying the Forchheimer / variable-porosity discretization}
\author{Porous Navier--Stokes solver --- verification note}
\date{\today}

\begin{document}
\maketitle

\begin{abstract}
This note documents a Method-of-Manufactured-Solutions (MMS) test whose
\emph{governing operator is identical} to the Cocquet benchmark --- the same
symmetric-gradient viscous term and the same nonlinear, porosity-dependent
Forchheimer--Ergun reaction --- but with a smooth, closed-form exact solution
and standard Dirichlet data. The Cocquet benchmark exhibits sub-optimal
$L^2$ convergence on coarse meshes; pure Navier--Stokes ($\alpha\equiv 1$,
$\sigma=0$) on the same geometry recovers optimal rates. The variable porosity
$\alpha$ \emph{combined with} the nonlinear drag is therefore the suspect, yet
this exact combination had never been exercised by a manufactured solution: the
existing MMS oracle only admitted a \emph{constant} reaction coefficient. The
test below closes that gap. If it converges optimally, the Cocquet
sub-optimality is physical/pre-asymptotic; if it degrades, a genuine
formulation or implementation defect has been isolated. The Cocquet experiment
itself is left untouched.
\end{abstract}

\section{Domain and porosity field}

The domain is the square $\Omega=[-\tfrac12,\tfrac12]^2$ with characteristic
length $L=1$. The porosity $\alpha:\Omega\to(0,1]$ is the smooth radial field
\texttt{SmoothRadialPorosity} used throughout the existing MMS suite. With
$r=\lvert\xx\rvert=\sqrt{x_1^2+x_2^2}$ and inner/outer radii $0<r_1<r_2<L/2$,
\begin{equation}
\alpha(\xx)=
\begin{cases}
\alpha_0, & r\le r_1,\\[2pt]
\alpha_\infty-\dfrac{\alpha_\infty-\alpha_0}{1+e^{\gamma(\eta)}}, & r_1<r<r_2,\\[8pt]
\alpha_\infty, & r\ge r_2,
\end{cases}
\qquad
\eta=\frac{r^2-r_1^2}{r_2^2-r_1^2},\quad
\gamma(\eta)=\frac{2\eta-1}{\eta(1-\eta)} .
\label{eq:alpha}
\end{equation}
The field is $C^\infty$ in the transition annulus and matches $C^1$ at
$r=r_1,r_2$; its gradient $\grad\alpha$ and Laplacian $\Delta\alpha$ are
computed analytically (closed-form chain rule, \texttt{src/models/porosity.jl}).
In the test we use $r_1=0.2$, $r_2=0.4$, $\alpha_\infty=1$, and sweep
$\alpha_0\in\{1.0,\,0.5,\,0.05\}$. The case $\alpha_0=1$ gives
$\alpha\equiv1$, $\grad\alpha\equiv\bm 0$ and $\sigma\equiv0$ --- a pure
Navier--Stokes control. The cases $\alpha_0<1$ activate both the porosity
gradient and the Forchheimer drag.

\section{Exact solution}

The manufactured velocity and pressure are
\begin{align}
\uu_{\mathrm{ex}}(\xx)
  &= U\,\alpha_0\,\frac{1}{\alpha(\xx)}\,\bm S(\xx),
  &\bm S(\xx)&=\begin{pmatrix}\sin(\pi x_1)\sin(\pi x_2)\\[2pt]
                              \cos(\pi x_1)\cos(\pi x_2)\end{pmatrix},
  \label{eq:uex}\\[4pt]
p_{\mathrm{ex}}(\xx) &= P_c\,\cos(\pi x_1)\sin(\pi x_2),
  \label{eq:pex}
\end{align}
with amplitude $U=1$. The $1/\alpha$ weighting couples the velocity field to
the porosity, so that $\grad\uu_{\mathrm{ex}}$ carries an explicit
$\grad\alpha$ contribution and the manufactured problem genuinely exercises the
variable-media terms. Exact $\grad\uu_{\mathrm{ex}}$ and
$\Delta\uu_{\mathrm{ex}}$ follow from \eqref{eq:uex} by differentiating the
product $\alpha_0\,\alpha^{-1}\bm S$ (see \texttt{UExFunc} in
\texttt{src/problems/mms\_paper\_2d.jl}). The pressure scale is
\begin{equation}
P_c=(1+\mathrm{Re}+\mathrm{Da})\,\frac{U\nu}{L},
\end{equation}
which for the Forchheimer law reduces to $P_c=(1+\mathrm{Re})\,U\nu/L$ since no
single constant $\sigma$ (hence no Darcy number $\mathrm{Da}$) is defined.

\section{Governing equations (the Cocquet formulation)}

The steady, dimensional Darcy--Brinkman--Forchheimer momentum and mass balance
solved by the code are
\begin{align}
\alpha\,(\uu\!\cdot\!\grad)\uu
  \;+\;\alpha\,\grad p
  \;+\;\sigma(\alpha,\uu)\,\uu
  \;-\;\dive\!\big(2\nu\,\alpha\,\eps(\uu)\big)
  &= \ff,
  \label{eq:mom}\\[4pt]
\dive(\alpha\,\uu)\;+\;\epsilon_{\mathrm p}\,p &= g,
  \label{eq:mass}
\end{align}
on $\Omega$, where $\nu$ is the kinematic viscosity, $\epsilon_{\mathrm p}\ll1$
a pressure-penalty regularization, and $\ff$, $g$ the manufactured sources
constructed in \S\ref{sec:forcing}. The Reynolds number sets the viscosity,
$\mathrm{Re}=UL/\nu \Rightarrow \nu=UL/\mathrm{Re}$.

\subsection{Viscous term: symmetric gradient}
\label{sec:visc}
The Cocquet configuration uses the \emph{symmetric-gradient} viscous operator
(\texttt{viscous\_operator\_type = "SymmetricGradient"}), i.e. the strain-rate
tensor $\eps(\uu)=\tfrac12(\grad\uu+\grad\uu^{\mathsf T})$ \emph{without} the
deviatoric projection. Expanding the divergence with the product rule,
\begin{equation}
\dive\!\big(2\nu\,\alpha\,\eps(\uu)\big)
   = 2\nu\,\eps(\uu)\,\grad\alpha
     \;+\;\nu\,\alpha\,\Delta\uu
     \;+\;\nu\,\alpha\,\grad(\dive\uu),
  \label{eq:visc-strong}
\end{equation}
using $\dive\eps(\uu)=\tfrac12\Delta\uu+\tfrac12\grad(\dive\uu)$. All three
terms are retained in the strong residual used for the stabilization and in the
manufactured forcing. (The canonical deviatoric-symmetric operator would
instead drop the $\grad(\dive\uu)$ contribution in 2D; the test deliberately
matches Cocquet by using the symmetric-gradient form.)

\subsection{Reaction term: Forchheimer--Ergun}
\label{sec:reac}
The drag is the nonlinear, porosity-dependent Forchheimer--Ergun law
\begin{equation}
\sigma(\alpha,\uu)
   = \sigma_{\mathrm{lin}}\left(\frac{1-\alpha}{\alpha}\right)^{\!2}
   + \sigma_{\mathrm{nl}}\left(\frac{1-\alpha}{\alpha}\right)\lvert\uu\rvert_\star ,
  \label{eq:sigma}
\end{equation}
with the test coefficients $\sigma_{\mathrm{lin}}=0.30$,
$\sigma_{\mathrm{nl}}=1.75$ (identical to Cocquet). The speed
$\lvert\uu\rvert_\star$ is the regularized magnitude returned by
\texttt{effective\_speed}:
\begin{equation}
\lvert\uu\rvert_\star=\sqrt{\uu\!\cdot\!\uu+u_{\mathrm{floor}}^2},
\qquad
u_{\mathrm{floor}} = u_{\mathrm{base}}
   + w_h\,\frac{c_1\nu}{c_2 h + \epsilon_{\mathrm{floor}}}.
  \label{eq:floor}
\end{equation}
The MMS configuration fixes the mesh-coupling weight $w_h=0$, so
$u_{\mathrm{floor}}=u_{\mathrm{base}}$ is a \emph{mesh-independent} constant.
This is essential: it lets the forcing be evaluated by a point-wise oracle that
has no element size $h$, while remaining \emph{exactly} consistent with the
assembled reaction term. The forcing oracle calls the very same
\texttt{sigma}/\texttt{effective\_speed} routines as the assembly, so
\eqref{eq:sigma}--\eqref{eq:floor} are reproduced bit-for-bit. (The
implementation guards $w_h=0$ and errors otherwise.)

Because $\sigma$ depends on $\alpha(\xx)$ through $(1-\alpha)/\alpha$ and on
$\lvert\uu\rvert$, the reaction is simultaneously \emph{spatially varying} and
\emph{nonlinear in $\uu$} --- the precise feature absent from every previous
MMS run, all of which used a constant $\sigma$.

\section{Manufactured forcing}
\label{sec:forcing}
The sources $\ff$, $g$ are the strong residuals of
\eqref{eq:mom}--\eqref{eq:mass} evaluated at the exact fields
\eqref{eq:uex}--\eqref{eq:pex}:
\begin{align}
\ff(\xx) &=
   \alpha\,(\uu_{\mathrm{ex}}\!\cdot\!\grad)\uu_{\mathrm{ex}}
   + \alpha\,\grad p_{\mathrm{ex}}
   + \sigma(\alpha,\uu_{\mathrm{ex}})\,\uu_{\mathrm{ex}}
   - \Big[\,2\nu\,\eps(\uu_{\mathrm{ex}})\,\grad\alpha
           + \nu\,\alpha\,\Delta\uu_{\mathrm{ex}}
           + \nu\,\alpha\,\grad(\dive\uu_{\mathrm{ex}})\,\Big],
  \label{eq:f}\\[4pt]
g(\xx) &= \epsilon_{\mathrm p}\,p_{\mathrm{ex}}
   + \alpha\,\dive\uu_{\mathrm{ex}}
   + \grad\alpha\!\cdot\!\uu_{\mathrm{ex}} .
  \label{eq:g}
\end{align}
Every term is computed analytically, including
$\grad(\dive\uu_{\mathrm{ex}})$: since the base shape
$\bm S$ is divergence-free ($\dive\bm S\equiv0$), one has
$\dive\uu_{\mathrm{ex}}=-U\alpha_0(\bm S\!\cdot\!\grad\alpha)/\alpha^2$, whose
gradient follows in closed form from the exact porosity gradient $\grad\alpha$
and Hessian $\grad^2\alpha$ (\texttt{hess\_alpha}, \texttt{grad\_div\_u\_ex}).
This replaces an earlier finite-difference evaluation, removing all step-size
error. With $\ff$, $g$ so defined, the strong residual vanishes
identically at $\uu_{\mathrm{ex}},p_{\mathrm{ex}}$, and the stabilized scheme
is consistent by construction --- exactly as in the constant-$\sigma$ MMS, now
extended to \eqref{eq:sigma}.

\section{Boundary conditions}

The manufactured test imposes the exact velocity as a Dirichlet condition on
the \emph{entire} boundary, $\uu=\uu_{\mathrm{ex}}$ on $\partial\Omega$; the
pressure is left free and pinned only through the penalty
$\epsilon_{\mathrm p}p$ in \eqref{eq:mass}. This differs from the Cocquet
benchmark, which prescribes a parabolic inlet, no-slip walls, and a
traction-free Neumann outlet. The boundary \emph{data} therefore differ, but
the differential operator, viscous form, reaction law and coefficients,
stabilization (ASGS), and element pair are identical. Pure Navier--Stokes
runs (\S1, $\alpha_0=1$) already confirmed that the Cocquet boundary set does
not, by itself, spoil convergence, so isolating the operator under clean
Dirichlet data is the sharper diagnostic.

\section{Discretization and expected rates}

Taylor--Hood $\mathbb{P}_2/\mathbb{P}_1$ elements on simplicial (TRI) meshes,
ASGS stabilization, matching Cocquet. The porosity is interpolated to the
velocity space order so it introduces no extra modeling error relative to the
$\mathbb{P}_2$ velocity. The quadrature degree includes the Forchheimer bump
$4k_v+\lfloor k_v/2\rfloor$ to integrate the non-polynomial
$\lvert\uu\rvert$ factor in \eqref{eq:sigma}.

For a smooth solution the expected asymptotic $L^2$ rates are
\begin{equation}
\lVert\uu-\uu_h\rVert_{L^2}=O(h^{k_v+1})=O(h^3),
\qquad
\lVert p-p_h\rVert_{L^2}=O(h^{k_p+1})=O(h^2),
\end{equation}
with semi-$H^1$ rates one order lower. The sweep runs
$N\in\{10,20,40,80\}$ and $\mathrm{Re}\in\{1,100,500\}$ (the last matching
Cocquet), recording per-mesh errors and fitted slopes.

\paragraph{Reading the result.}
\begin{itemize}
\item \textbf{Optimal rates} (slopes $\to 3$ velocity, $\to 2$ pressure):
      the Forchheimer / variable-$\alpha$ discretization is consistent and
      accurate. The Cocquet sub-optimality is then attributable to physics
      (boundary/corner regularity, high-$\mathrm{Re}$ pre-asymptotic range, or
      a not-yet-converged fine reference), not to the formulation.
\item \textbf{Degraded rates:} a real consistency or implementation defect in
      the variable-porosity nonlinear-reaction path has been isolated;
      candidates to examine next are the velocity-floor consistency near
      stagnation, the symmetric-gradient operator under variable $\alpha$, the
      omitted porosity-gradient correction $C_\alpha=\alpha+(h/|k_0|)|\grad\alpha|$
      in $\tau$, and the reaction adjoint sign/weighting.
\end{itemize}

\section{Results}
\label{sec:results}

The sweep ($N\in\{10,20,40,80\}$, TRI, $\mathbb P_2/\mathbb P_1$, ASGS) gives
the fitted $L^2$ slopes below. The case $\alpha_0=1$ is the pure-NS control.

\begin{center}
\begin{tabular}{lccc}
\toprule
$(\alpha_0,\mathrm{Re})$ & rate $u_{L^2}$ & rate $p_{L^2}$ & rate $u_{H^1}$\\
\midrule
$(1.0,\,1)$    & 3.31 & 2.02 & 2.17\\
$(1.0,\,100)$  & 3.24 & 2.02 & 2.13\\
$(1.0,\,500)$  & 2.94 & 2.03 & 1.90\\
$(0.5,\,1)$    & 2.93 & 2.68 & 1.86\\
$(0.5,\,100)$  & 2.70 & 2.03 & 1.70\\
$(0.5,\,500)$  & 2.80 & 2.03 & 1.76\\
$(0.05,\,1)$   & 2.87 & 2.69 & 1.79\\
$(0.05,\,100)$ & 2.35 & 1.71 & 1.38\\
$(0.05,\,500)$ & 2.39 & 1.99 & 1.40\\
\bottomrule
\end{tabular}
\end{center}

The per-refinement velocity slopes \emph{increase} toward the optimal $3$ as
$h\to0$ in every case --- e.g. the stiffest run $(\alpha_0=0.05,\mathrm{Re}=500)$
gives $1.97\to2.54\to2.60$, and the control $(\alpha_0=1,\mathrm{Re}=500)$ sits
flat at $\approx 2.9$. There is a mild pre-asymptotic erosion on coarse meshes
that grows with the drag strength (small $\alpha_0$) and the Reynolds number,
but \emph{no fixed sub-optimal order ceiling}: the scheme is consistent and
optimally convergent for the nonlinear, variable-porosity formulation.

\paragraph{Conclusion.}
The Cocquet \emph{formulation} is correct. An identical operator on a smooth
solution recovers $O(h^3)$ velocity / $O(h^2)$ pressure, whereas the Cocquet
benchmark itself plateaued at $\approx0.1$--$0.5$. The benchmark's
sub-optimality is therefore \emph{not} a formulation/implementation defect in
the discretized PDE; it had to be something specific to the benchmark
\emph{run} that this manufactured test does not share. That cause was
subsequently identified (see \texttt{docs/cocquet\_convergence\_analysis.md},
Phase~6): the velocity \emph{regularization floor} used for the stabilization
parameters $\tau$ carries a mesh-dependent diffusive term
$\propto \nu/h$ that was also being fed into the \emph{physical} Forchheimer
drag $\beta(\alpha)\lvert\uu\rvert\uu$. Growing under refinement, it made every
mesh solve a different effective drag law and corrupted the reference solution,
collapsing the measured rates. Removing that term from the reaction speed
(while keeping it in $\tau$) restores velocity $H^1$ slopes of $\approx 2$ and
$L^2$ slopes of $2.3$--$2.8$ for the benchmark, consistent with the present
manufactured study and with the source paper's reported $O(h^2)$. This MMS,
crucially, is what proved the discretization itself was sound and pointed the
search away from the formulation.

\section*{Implementation note}
The manufactured-solution oracle in \texttt{src/problems/mms\_paper\_2d.jl} is
\emph{shared} between the standard MMS suite and this Cocquet-form test. It
dispatches on the configured operators, so the two formulations differ in
exactly two terms --- the viscous operator (symmetric-gradient vs.
deviatoric-symmetric) and the reaction law (constant $\sigma$ vs. nonlinear
Forchheimer--Ergun) --- and the single oracle reproduces both forcings exactly.
This isolates precisely what changes between the formulations.

\section*{Reproduce}
\begin{verbatim}
cd test/extended/CocquetFormMMS
julia --project=../../.. run_test.jl cocquet_form_mms.json
# rates stored as HDF5 attributes rate_u_l2, rate_p_l2, ... in
# results/cocquet_form_mms.h5
\end{verbatim}

\end{document}
